\documentclass[a4paper]{article}
\usepackage[T2A]{fontenc}
\usepackage[utf8]{inputenc}
\usepackage[russian]{babel}
\usepackage{geometry}
\usepackage{setspace} % Для управления интервалами
\usepackage{myrussian}

\geometry{
  left=3cm,
  right=1cm,
  top=2cm,
  bottom=2cm
}

\begin{document}
\textbf{4.} Положим \( \sigma _1, \sigma _2 \) --- образующие,
соответствующие перестановке, в которой нить идущая слева направо
лежит сверху (рис. \ref{fig:sigma1 and 2}).

\epsimage{sigma1}{0.14}{fig:sigma1 and 2}{\( \sigma _1, \sigma _2\)}

Чтобы доказать, что \(\sigma _1 \neq \sigma _2 \)
покажем, что \( \sigma _1 \sigma _2^{-1} \neq 1 \). Если бы это было
так, то соответствующие диаграммы были бы изотопны. Тогда должны быть
изотопны и их замыкания. Диаграмма соответствующая единице --- три
прямых отрезка, ее замыкание --- тривиальное зацепление из трех
компонент. На рисунке \ref{fig:sigma1 sigma 2 -1} показано замыкание
диаграммы \( \sigma _1
\sigma _2^{-1} \). На рис. \ref{fig:sigma12iso} показано изотопное
зацепление, которое
изотопно тривиальному узлу. В частности оно не изотопно зацеплению из
трех компонент.

\epsimage{sigma12}{0.20}{fig:sigma1 sigma 2 -1}{Замыкание \( \sigma
_1 \sigma _2^{-1}\)}

\epsimage{sigma12iso}{0.20}{fig:sigma12iso}{Сглаженное замыкание \( \sigma
_1 \sigma _2^{-1}\)}
\end{document}
